% Controlled Counterfactual Perturbation -- drop-in section.
% Uses only what the main paper already loads (\[ \], \lVert, \mathbb, \mathrm, \texttt, tabular).
% All values are filled from run 20260924-153459 (commit 758b7c9); the artefact that fills
% each table is named in a comment above it, and \RESULT{...} is kept for future runs.
% Notation follows the transcoder section of the main paper.

\providecommand{\RESULT}[1]{\textbf{[#1]}}

\section{Controlled Counterfactual Perturbation}
\label{sec:counterfactual}

The preceding sections nominate a feature by how unusually it activates, trace
a circuit upstream of it by attribution, and test that circuit by ablating it
and measuring the action. Each step is selection-based: positives are the
observations that activated the feature, negatives are observations judged
related, and the model is never shown an input whose only difference from
another is known. As a result the pipeline can establish that a traced set is
\emph{load-bearing}, since removing it moves the action, but not that it
\emph{carries} the content the target feature is claimed to encode. A set of
active, high-variance features at the right layers would also move the action
when removed.

This section supplies the missing kind of evidence: an intervention on the
\emph{input} whose cause is fixed by construction. We freeze one moment of a
simulated scene, recolour exactly one object, and push the original and
recoloured frames through the unmodified policy with the same sampling noise,
so that any change in the sparse code is attributable to that object's pixels.
A visually identical object that the task does not refer to serves as a matched
placebo, a prompt swap over identical pixels tests whether the response follows
the language, and the same controlled cause is then used to audit the circuit
tracer of Section~B: parents that genuinely carry a feature's information about
the object should respond to the object's recolour more than random features
at their layers do.

Three findings result. Transcoder features respond to the referent's recolour
several times more than to the identical non-referent's, and the policy's
action follows suit; the prompt swap available in this task suite does not move
the behavioural referent, so the question of language versus position is not
decidable on this scene; and for a strongly object-selective target, the
tracer's parents are exactly as object-selective as their layers and no more,
while passing a response-only test of the kind Section~C uses. The last of
these is the result the section was designed to be able to produce.

\subsection{Hypotheses}
\label{sec:hypotheses}

\paragraph{H1, object selectivity.}
Transcoder features respond to a controlled appearance change of the object
the task designates more than to a matched appearance change of a visually
identical object it does not, once perturbation size is accounted for.
Falsified if the two responses are comparable after that adjustment, which
would indicate the features track generic pixel change. The two objects
necessarily differ in screen position as well as in whether the prompt refers
to them; H1 alone cannot separate these.

\paragraph{H2, referent versus position.}
The selectivity of H1 follows the object the language refers to. Supported if
swapping the prompt to refer to the other object, over pixel-identical images,
inverts the selectivity; falsified if selectivity is unchanged or strengthens
for the original object; partial if it weakens without inverting. All three
readings presuppose that the swap moved the referent, which is checked on the
action rather than assumed.

\paragraph{H3, circuit corroboration.}
The parents of a circuit traced upstream of a target are more selective for the
perturbed object than random features matched in layer and flow time and drawn
from those the scene exercises. A single circuit can only ever produce a single
anecdote, and H3 is a claim about the tracer rather than about one of its
outputs, so the decision quantity is the fraction $k/N$ of nominated targets
whose circuits clear that bar, over targets pooled across tasks. Supported if
the exact interval on $k/N$ excludes the audit's own false-positive rate;
falsified if its upper end falls below a half, so that the tracer fails more
often than it works; inconclusive otherwise. Per target the bar is a
selectivity enrichment above one at the head of the influence ranking. Whether
the parents are also enriched for a manipulation of a different kind qualifies
a supported verdict but does not overturn it, and the target node is excluded
from the test because it was nominated for responding.

A null result here is a claim about absence, and absence is only readable
against what the instrument could have seen. The audit's sensitivity is
therefore measured rather than assumed (Sec.~\ref{sec:calibration}), and the
size of $N$ is chosen so that a null result excludes something
(Sec.~\ref{sec:rate}).

H1 and H2 are claims about the model. H3 is a claim about the method: it uses a
known cause to ask whether attribution-based tracing recovers the features that
carry that cause.

\subsection{Method}

\subsubsection{Paired same-state counterfactual}

Fix a task and a seed and advance the simulator to a state $s$. Let
$R(s,\theta)$ be the rendered observation under scene appearance $\theta$, with
$\theta_0$ the unmodified scene and $\theta_c^{(d)}$ the scene with object $c$
recoloured at dose $d$. A forward pass yields, at every instrumented layer
$\ell$ and flow time $\tau$, a sparse code $z_{\ell,\tau}\in\mathbb{R}^{B\times
T\times d_{\mathrm{features}}}$ with $B=1$ and $T=50$ action tokens. We reduce
over batch and tokens by a maximum,
\begin{equation}
\tilde{z}_{\ell,\tau,i} = \max_{b,t} z_{\ell,\tau,b,t,i},
\label{eq:reduce}
\end{equation}
because the question is whether a feature fired anywhere in the chunk, not how
broadly. With $\epsilon$ a single flow-matching noise sample shared by both
passes, the signed per-feature difference is
\begin{equation}
\delta_{\ell,\tau,i}(c,d)
=
\tilde{z}_{\ell,\tau,i}\!\left(R(s,\theta_c^{(d)}),\epsilon\right)
-
\tilde{z}_{\ell,\tau,i}\!\left(R(s,\theta_0),\epsilon\right),
\label{eq:delta}
\end{equation}
recorded in full for every layer, flow time and condition. Two summaries are
used. The \emph{layer response} is the norm over features averaged over flow
times,
\begin{equation}
D_{\ell}(c,d) = \frac{1}{|\mathcal{T}|}\sum_{\tau\in\mathcal{T}}
\left\lVert \delta_{\ell,\tau,\cdot}(c,d) \right\rVert_2 ,
\label{eq:layer-response}
\end{equation}
and drives the headline selectivity; the \emph{per-feature response}
$|\delta_{\ell,\tau,i}|$ drives the nomination of a trace target and the
circuit audit.

Three properties make Eq.~\ref{eq:delta} an intervention rather than a
comparison. The state is frozen: physics is never stepped between renders, so
pose, camera and lighting are identical, and after the recolours at a state are
reverted the baseline is re-rendered and required to match exactly. The noise
is shared and seeded: a generator seeded by (seed, state, draw) produces one
sample per replicate that both passes receive, so runs are reproducible bit for
bit. The loop is open: between states the simulator is advanced by executing
the first $k$ actions of the policy's own unperturbed plan, and no measurement
is taken along a diverging closed-loop rollout.

A \emph{cell} is one combination of state, noise draw and dose under one
prompt. Target and placebo are measured in every cell; the null and positive
controls, which have no dose, once per (state, draw). Colour is the
manipulation because it changes no geometry, mass, contact or position, and in
this suite it is also load-bearing: every prompt names ``the black bowl'', so
recolouring a bowl removes it from what the noun denotes.

\subsubsection{Matched placebo and footprint adjustment}
\label{sec:placebo}

The suite instantiates two copies of the same mesh, and the task names one.
Perturbing the other gives a control matched on mesh, material, colour and
size, differing only in table position and in whether the prompt refers to it.
It is the \emph{non-referent}, not an irrelevant object: the policy must
discriminate against it to solve the task, so H1 asks whether features weight
the referent above the non-referent, not whether they ignore it.

Writing $D(c)$ for the layer response pooled over all $18$ layers, all $n$ flow
times and all cells under the task prompt, the raw selectivity is
$\mathrm{Sel}=D(\mathrm{target})/D(\mathrm{placebo})$. Because the two
instances occupy different screen positions, their footprints differ, and we
also report
\begin{equation}
\mathrm{Sel}_{\mathrm{adj}}
=
\frac{D(\mathrm{target})}{D(\mathrm{placebo})}
\bigg/
\frac{S(\mathrm{target})}{S(\mathrm{placebo})},
\qquad
S(c)\in\{\,|\mathcal{P}(c)|,\ \lVert\Delta I(c)\rVert_2\,\},
\label{eq:sel-px}
\end{equation}
with $\mathcal{P}$ the changed pixels and $\Delta I$ the image difference,
each computed over every camera the policy sees, since a static camera's view
of a static bowl does not change between states while the wrist camera's does.
The image norm is the primary normaliser: unlike the pixel count it grows with
dose. Eq.~\ref{eq:sel-px} is exact only under a linear response, which we test
as a dose elasticity $\eta = \Delta\log D/\Delta\log S$ between consecutive
doses; $\eta=1$ is linear.

Uncertainty is reported at two levels. Cells give the mean, standard deviation,
range and a percentile bootstrap of the pooled ratio; but doses within a
(state, draw) share a render and a noise sample, so the range over (state,
draw) clusters is reported as well. The pre-registered H1 rule uses the minimum
over cells, which assumes no independence. Because $D$ is a norm over $16384$
features, the layers with the largest activations dominate it; we therefore
also report the ratio on the relative layer response
$\lVert\delta\rVert_2/\lVert\tilde{z}\rVert_2$ and the geometric mean over
layers of per-layer selectivities, and call the conclusion robust only if all
three lie on the same side of one.

\subsubsection{Controls}
\label{sec:controls}

\emph{Null.} The same baseline is pushed through the policy twice with the same
noise; the resulting $\delta$ is the nondeterminism floor, and a run is valid
only if its largest layer response is at most $5\%$ of the placebo's.
\emph{Positive control.} The prompt swap is a manipulation known to change
behaviour; its latent response, available from the two baselines already
computed, gives the scale against which a recolour response is read.
\emph{Liveness.} Before measurement the target is recoloured to saturation and
the pixel change required to be nonzero, then reverted and the baseline
required to be restored exactly, which a stale render cache would fail.
\emph{Dose.} Two doses along a linear blend, with $\eta$ and monotonicity
reported per condition. \emph{Guards.} Tensor shapes are checked on every pass
against the first; the policy loader's report is checked for missing
parameters, since it otherwise returns random weights with a warning.

\subsubsection{Referent versus position}
\label{sec:grounding}

Task relevance and screen position are confounded between the two bowls, and
the suite defines its targets by position, so they cannot be separated by
moving objects. They can be separated by language: the suite's sibling tasks
place the same object in different regions, and the sibling whose target
occupies the placebo's region supplies a prompt that refers to it. Running both
prompts over identical images with shared noise leaves language as the only
difference. Selectivity is computed within each prompt against that prompt's
own baseline, so the prompt's large direct effect on the activations cancels.

Two checks gate the comparison. Prompt grounding,
\begin{equation}
g = \frac{\lVert a(\mathrm{prompt}_{\mathrm{alt}}) - a(\mathrm{prompt}_{\mathrm{task}})\rVert_2}
         {\lVert a(\mathrm{prompt}_{\mathrm{task}})\rVert_2},
\label{eq:grounding}
\end{equation}
shows the policy processed the new words; below $0.01$ the swap is inert. But
processing is not re-grounding. Because the prompt names ``the black bowl'',
recolouring the bowl the policy is going for moves the action far more than
recolouring the other, so the action's own sensitivity reveals the behavioural
referent. With $A(c)$ the relative action change under recolour of $c$,
\begin{equation}
\mathrm{anchor}(\mathrm{prompt}) = \frac{A(\mathrm{target})}{A(\mathrm{placebo})}
\label{eq:anchor}
\end{equation}
must exceed one under the task prompt and fall below one under the sibling
prompt for the referent to have moved. If it does not, the swap moved words and
H2 is untestable on this scene with this prompt, whatever $g$ says.

\subsubsection{Nominating a trace target}
\label{sec:zscore}

The probe nominates the feature whose circuit is traced. Raw selectivity cannot
choose across depth, because placebo responses differ by an order of magnitude
between layers, so candidates are ranked by a layer-standardised response
\begin{equation}
z^{\mathrm{std}}_{\ell,i}
=
\frac{\overline{|\delta|}_{\ell,i}-\mu_\ell}{\sigma_\ell},
\label{eq:zstd}
\end{equation}
where $\overline{|\delta|}_{\ell,i}$ is the feature's mean response over the
cells in which it ranked among the $K$ largest responders at its layer and flow
time, and $\mu_\ell,\sigma_\ell$ summarise that quantity over the layer's
recorded features. Three filters precede ranking: the feature must rank among
the top $K$ in at least $25\%$ of the layer's cells, so a single large
observation cannot nominate it; it must sit at layer $\ell\ge4$, so there are
parents to find; and its exact placebo response at its peak flow time must be
at most half its target response. The placebo is a filter, not a denominator,
so a placebo response of zero is admissible.

\subsubsection{Auditing the traced circuit}
\label{sec:validation-method}

Section~B traces an edge from feature $i$ at layer $\ell$ to feature $j$ at a
later layer as the attribution
\begin{equation}
\mathrm{edge}(i\rightarrow j)=\sum_{t=1}^{T} z_{\ell,i,t}\,
\frac{\partial a_{\ell',j,p^{\star}}}{\partial z_{\ell,i,t}},
\label{eq:edge}
\end{equation}
the first-order estimate of how much of $j$'s pre-activation at position
$p^{\star}$ owes to $i$ being active; the derivative is taken in the
replacement model and the source is summed over its token positions. Two
things follow. The edge is a local, first-order quantity that need not carry
semantic content, which is what H3 tests; and it favours parents that are
active and high-variance on the tracing observations, which is the rival
explanation H3 must rule out.

Let $\mathcal{C}$ be the parents of the traced circuit. Each is scored at its
own flow time as the mean $|\delta|$ over the task-prompt cells; because every
delta is recorded, every parent has a known response. A parent is
\emph{exercised} if it is active in some baseline or moves under the target
perturbation in some cell; unexercised parents are excluded and counted. The
null is matched: for each exercised parent a random exercised feature is drawn
at the same layer and flow time, the same draw is evaluated under every
condition, and with $M$ draws the response enrichment is
\begin{equation}
\mathrm{Enrich}(\mathcal{C},c)=\frac{\bar{\Delta}_{\mathcal{C}}(c)}
{\frac{1}{M}\sum_m \bar{\Delta}^{(m)}(c)},
\qquad
p=\frac{1+|\{m:\bar{\Delta}^{(m)}(c)\ge\bar{\Delta}_{\mathcal{C}}(c)\}|}{M+1}.
\label{eq:enrich}
\end{equation}

Response enrichment cannot decide H3: a generically responsive set clears it
for every condition. Generic responsiveness, however, scales the target and
placebo responses alike and leaves their ratio at the random set's, so the
deciding statistic is the ratio of ratios
\begin{equation}
\mathrm{SelEnrich}(\mathcal{C})=\frac{\mathrm{Sel}_{\mathcal{C}}}
{\frac{1}{M}\sum_m \mathrm{Sel}^{(m)}},
\qquad
\mathrm{Sel}_{\mathcal{C}}=\frac{\bar{\Delta}_{\mathcal{C}}(\mathrm{target})}
{\bar{\Delta}_{\mathcal{C}}(\mathrm{placebo})},
\label{eq:sel-enrich}
\end{equation}
with its own Monte-Carlo tail. A matched random set at these layers is itself
highly selective, because selectivity is a property of the layer as much as of
the feature, which is why $\mathrm{Sel}_{\mathcal{C}}$ alone is not evidence.
Specificity,
\begin{equation}
\mathrm{Spec}(\mathcal{C})=\frac{\mathrm{Enrich}(\mathcal{C},\mathrm{target})}
{\max_{c\in\mathcal{M}}\mathrm{Enrich}(\mathcal{C},c)},
\label{eq:specificity}
\end{equation}
over the manipulations $\mathcal{M}$ of a different kind (here the prompt swap;
the placebo is excluded as it is already the denominator above) qualifies a
supported verdict: above one the parents are object-specific, at or below one
they are as enriched for a language change, which is what a circuit that
integrates the referent from vision and language would also show. Specificity
alone therefore never falsifies.

At large $|\mathcal{C}|$ both means concentrate and any systematic difference
drives $p$ to its floor $1/(M+1)$ regardless of size; a selectivity enrichment
of $1.05$ over $1700$ nodes is significant and means nothing. The decision is
therefore taken where $p$ is informative: parents are ordered by traced
influence and H3 is decided on the top $10$, the tracer's strongest claim, with
the full circuit required to agree in direction. Enrichment is also reported
over the top $50$, top $200$ and all parents, since a ranking that carries
information shows a strong head that decays while a flat profile says it does
not, and the two call for different remedies. The nominee is chosen from the
probe's data and excluded from the test; the parents are chosen by the tracer
from dataset observations it never sees; nothing is both selected on and tested
with.

\subsubsection{Calibrating the audit}
\label{sec:calibration}

The audit of Eq.~\ref{eq:sel-enrich} has a negative control and no positive
one. A flat reading is therefore ambiguous between \emph{this circuit carries
nothing} and \emph{this test cannot see what it is looking for}, and the second
reading would void every negative verdict the audit has produced. The fix is to
run the audit on circuits whose content is known by construction.

Let $\mathcal{P}$ be the exercised features below the target's layer at its
flow time, each with the measured selectivity $\mathrm{sel}(f)=
\bar{\delta}_f(\mathrm{target})/\bar{\delta}_f(\mathrm{placebo})$ the probe
already recorded for it. For a contamination level $c\in[0,1]$ and a set size
$K$ matching the circuit under audit, build
\begin{equation}
\mathcal{S}(c)=\underbrace{\{f_1,\dots,f_{\lceil cK\rfloor}\}}
_{\text{sampled from the top of }\mathrm{sel}}
\;\cup\;
\underbrace{\{g_1,\dots,g_{K-\lceil cK\rfloor}\}}_{\text{uniform on }\mathcal{P}},
\label{eq:contamination}
\end{equation}
and run the audit's own statistic on it. At $c=0$ the set is random and the
audit must not call it; at $c=1$ it is maximally content-carrying and the audit
must. Repeating $R$ times per level, with the selective part drawn from a
widened head so that repeats differ, gives a detection rate per level, and the
smallest level detected at $80\%$ or better is the \emph{detection floor}
$c^{\star}$. A negative verdict then reads ``this circuit is less than
$c^{\star}$ content-carrying'' instead of ``this circuit was not enriched''.
Three outcomes are distinguished and named: an audit that misses $c=1$ is
blind and its negatives are void; one that calls $c=0$ more than $20\%$ of the
time over-calls and its positives are void; otherwise $c^{\star}$ is reported.
Selection uses the target condition's own measurements, so this calibrates the
instrument and is not itself a test of the model.

\subsubsection{From one circuit to a rate}
\label{sec:rate}

With $k$ of $N$ audited circuits corroborated, the interval on the underlying
rate is Clopper-Pearson, obtained by inverting the exact binomial tail. The
choice matters at the endpoint that decides a null H3: a Wald interval at $k=0$
has zero width and would certify a rate of exactly zero from a handful of
targets, whereas the exact one-sided bound is
\begin{equation}
p_{\max}=1-\alpha^{1/N},
\label{eq:cp-zero}
\end{equation}
which is $45\%$ at $N=5$, $26\%$ at $N=10$ and $14\%$ at $N=20$ with
$\alpha=0.05$. Five targets cannot falsify H3 however flat they are, because
the bound still straddles a half; twenty can. $N$ is therefore fixed at twenty
before the run, and targets are pooled across tasks so that the rate is a
property of the tracer rather than of one scene.

\subsection{Experimental setup}

% Filled from counterfactual_summary.json -> config, and observation_shapes.json.
\begin{table}[t]
\caption{Setup of the reported run. Cells per condition are states $\times$
draws $\times$ doses under one prompt; cells per layer multiply by the denoise
steps and are the denominator of the reproducibility filter.}
\label{tab:cf-config}
\begin{center}
\small
\begin{tabular}{p{1.35in}p{3.75in}}
\hline
Suite and tasks & \texttt{libero\_spatial}, all ten tasks. Every task in this suite instantiates two identical \texttt{akita\_black\_bowl} meshes, which is what makes a matched placebo possible; \texttt{libero\_object} and \texttt{libero\_goal} contain no duplicated object and admit none \\
Target, placebo & each task's own first object of interest, and the other instance of the same mesh, read from the task's problem definition \\
Sibling prompt & the task's phrasing with the referent replaced by the placebo's description \\
Seeds; states; draws; doses & $1000$, $1001$; $12$ per run (four plan actions apart); $2$; $0.5$, $1.0$ (blend toward red) \\
Cells per condition; per layer & $96$ per task; $960$ ($n=10$ denoise steps) \\
Task clusters & $10$, the unit the confidence interval resamples \\
Nomination & top-$K$ with $K=25$; $25\%$ consistency; $\ell\ge4$; exact selectivity $\ge2$ \\
Circuit audit & $N=20$ targets pooled across tasks; sources restricted to the previous layer; $M=2000$ matched draws; decision at the top $10$ by influence; $\alpha=0.05$ \\
Audit calibration & sets of $16$ from the exercised pool; $c\in\{0,\tfrac18,\tfrac14,\tfrac12,\tfrac34,1\}$; $R=20$ repeats \\
Shapes per pass & images $(1,3,360,360)\times2$; tokens $(1,200)$; state $(1,8)$; noise $(1,50,32)$; code $(1,50,16384)\times18\times10$; action $(50,7)$ \\
\hline
\end{tabular}
\end{center}
\end{table}

Table~\ref{tab:cf-config} records the design. Objects and the sibling prompt
are read from each task's own problem definition, so none of these choices is
the experimenter's. The task, not the cell, is the unit: cells within a task
share a scene, a camera geometry and a referent, so they are replicates of one
measurement rather than independent measurements, and the interval on H1 is a
cluster bootstrap that resamples whole tasks. Ten is the suite's ceiling rather
than a budget, and it is the reason H1 can carry an interval at all. Every stage records its configuration, the tensor shapes handed
to the policy, the git commit, library versions, device and a content hash of
the transcoder checkpoint; the language tokens carry the discretised
proprioceptive state and so differ between states while identical within a
cell.

\subsection{Results}

Two runs are reported and kept apart. The \emph{pilot} is a single task at a
single seed, and every number in Tables~\ref{tab:cf-main}
to~\ref{tab:cf-validation} comes from it. It established the design and
produced one falsification, and it cannot support an interval: four clusters
and one scene. The \emph{reported run} is the design of
Table~\ref{tab:cf-config}, ten tasks at two seeds with twenty audited
circuits, and fills Tables~\ref{tab:cf-pooled} and~\ref{tab:cf-rate}. The
audit calibration of Table~\ref{tab:cf-calibration} was run on the pilot's own
delta store and so applies to the pilot's negative directly.

\subsubsection{Object selectivity}

% Filled from decision_metrics.json -> h1 (pooled, spread, cluster_spread, bootstrap_ci_95,
% robustness, positive_control, dose_response, null_floor). Cells marked \RESULT depend on the
% perturbation size measured over both cameras.
\begin{table}[t]
\caption{Pilot, one task at one seed. Selectivity under the task prompt. $D$ is the pooled layer response of
Eq.~\ref{eq:layer-response}; $S$ the perturbation size over both cameras; the
null row is the nondeterminism floor. Eight cells; four (state, draw)
clusters.}
\label{tab:cf-main}
\begin{center}
\small
\begin{tabular}{lcccc}
\hline
Condition & $D$ & $S=|\mathcal{P}|$ & $S=\lVert\Delta I\rVert_2$ & Action rel.\ $\ell_2$ \\
\hline
Null control (max)            & $0$ & $0$ & $0$ & $0$ \\
Positive control, prompt swap & $19.7$ & $0$ & $0$ & $0.729$ \\
Target object                 & $10.11$ & $0.0148$ & $28.8$ & $0.538$ \\
Placebo object                & $1.42$  & $0.0049$ & $16.2$ & $0.032$ \\
\hline
Ratio, target over placebo    & $7.13$ & $3.03$ & $1.78$ & $17.0$ \\
\hline
\multicolumn{5}{p{5.1in}}{Raw $\mathrm{Sel}$ by cell: mean $8.56\pm5.95$, range $2.92$ to $15.7$, 95\% CI $[4.00,\,11.4]$; by cluster: $3.02$ to $14.4$} \\
\multicolumn{5}{p{5.1in}}{$\mathrm{Sel}_{\mathrm{adj}}$, Eq.~\ref{eq:sel-px}, $S=\lVert\Delta I\rVert_2$: pooled $4.01$; cells mean $6.55\pm5.69$, min $1.23$, max $13.2$; 95\% CI $[1.81,\,8.54]$; by cluster $1.27$ to $12.1$} \\
\multicolumn{5}{p{5.1in}}{$\mathrm{Sel}_{\mathrm{adj}}$, $S=|\mathcal{P}|$: pooled $2.36$; cells min $0.61$, max $12.8$; 95\% CI $[0.92,\,6.81]$} \\
\multicolumn{5}{p{5.1in}}{Robustness: per-layer geometric mean $7.33$ [$5.32$, $11.3$], $18$ of $18$ layers above 1; relative-$\ell_2$ ratio $7.07$; all three agree in direction} \\
\multicolumn{5}{p{5.1in}}{Positive control: the referent recolour is $0.51$ of the prompt swap's latent response, the placebo $0.07$} \\
\multicolumn{5}{p{5.1in}}{Dose: target $D$ $9.49\rightarrow10.7$, placebo $1.25\rightarrow1.59$, both monotone; $\eta(\mathrm{target})=0.62$, $\eta(\mathrm{placebo})=1.21$} \\
\hline
\end{tabular}
\end{center}
\end{table}

% Filled from layer_selectivity.csv, prompt = task.
\begin{table}[t]
\caption{Per-layer response $D_\ell$ under the task prompt. Selectivity is
comparable within a layer; across layers the layer-standardised score of
Eq.~\ref{eq:zstd} is used.}
\label{tab:cf-layers}
\begin{center}
\small
\begin{tabular}{cccc|cccc}
\hline
Layer & Target & Placebo & Sel. & Layer & Target & Placebo & Sel. \\
\hline
0 & $28.00$ & $2.470$ & $11.3$ & 9  & $3.444$ & $0.511$ & $6.7$ \\
1 & $17.80$ & $1.632$ & $10.9$ & 10 & $4.328$ & $0.615$ & $7.0$ \\
2 & $9.686$ & $0.900$ & $10.8$ & 11 & $3.652$ & $0.549$ & $6.6$ \\
3 & $2.427$ & $0.262$ & $9.3$  & 12 & $5.949$ & $0.812$ & $7.3$ \\
4 & $1.069$ & $0.141$ & $7.6$  & 13 & $9.110$ & $1.126$ & $8.1$ \\
5 & $1.205$ & $0.167$ & $7.2$  & 14 & $10.17$ & $1.882$ & $5.4$ \\
6 & $1.133$ & $0.189$ & $6.0$  & 15 & $13.79$ & $2.592$ & $5.3$ \\
7 & $3.046$ & $0.432$ & $7.0$  & 16 & $18.66$ & $2.993$ & $6.2$ \\
8 & $2.556$ & $0.387$ & $6.6$  & 17 & $46.02$ & $7.857$ & $5.9$ \\
\hline
\end{tabular}
\end{center}
\end{table}

% Filled from decision_metrics.json -> h1.cells (D, S, Sel) and measurements.csv (A).
\begin{table}[t]
\caption{Pilot, every cell under the task prompt. $D$ is the layer response of
Eq.~\ref{eq:layer-response} pooled over layers and flow times; $S$ the
image-difference norm over both cameras; $\mathrm{Sel}_{\mathrm{adj}}$ divides
the $D$ ratio by the $S$ ratio; $A$ the relative change in the predicted action
chunk. State 1 is five plan actions after state 0. The table is shown in full
because the gap between its two states is the reason the reported run samples
twelve: the reported run has $960$ such cells and they are released as a table
rather than printed.}
\label{tab:cf-cells}
\begin{center}
\footnotesize
\begin{tabular}{cccccccccc}
\hline
State & Draw & Dose & $D_{\mathrm{target}}$ & $D_{\mathrm{placebo}}$ & $S_{\mathrm{t}}/S_{\mathrm{p}}$ & $\mathrm{Sel}$ & $\mathrm{Sel}_{\mathrm{adj}}$ & $A_{\mathrm{target}}$ & $A_{\mathrm{placebo}}$ \\
\hline
0 & 0 & 0.5 & $4.67$  & $1.55$ & $2.37$ & $3.02$  & $1.27$  & $0.129$ & $0.030$ \\
0 & 0 & 1.0 & $6.65$  & $1.95$ & $2.37$ & $3.41$  & $1.44$  & $0.228$ & $0.046$ \\
0 & 1 & 0.5 & $4.65$  & $1.59$ & $2.37$ & $2.92$  & $1.23$  & $0.119$ & $0.031$ \\
0 & 1 & 1.0 & $6.15$  & $1.99$ & $2.37$ & $3.10$  & $1.31$  & $0.205$ & $0.049$ \\
1 & 0 & 0.5 & $14.30$ & $0.91$ & $1.19$ & $15.71$ & $13.22$ & $0.888$ & $0.021$ \\
1 & 0 & 1.0 & $15.10$ & $1.13$ & $1.19$ & $13.38$ & $11.25$ & $0.906$ & $0.024$ \\
1 & 1 & 0.5 & $14.33$ & $0.94$ & $1.19$ & $15.20$ & $12.79$ & $0.906$ & $0.025$ \\
1 & 1 & 1.0 & $15.05$ & $1.28$ & $1.19$ & $11.75$ & $9.88$  & $0.924$ & $0.027$ \\
\hline
\end{tabular}
\end{center}
\end{table}

The null floor is exactly zero in every cell: the forward pass is
bit-deterministic on this hardware, so every recorded difference is signal.
Recolouring the referent bowl moves the pooled layer response by $10.1$ and the
identical non-referent by $1.42$, a raw selectivity of $7.1$
(Table~\ref{tab:cf-main}). The ratio exceeds one in every one of the eight
cells (Table~\ref{tab:cf-cells}), with a minimum of $2.9$, and it is
heterogeneous across states: near $3$ at the initial state and near $14$ five
plan actions later. The two noise draws at a state agree closely, so the
heterogeneity is the state, not the sampling.

The footprint adjustment explains the heterogeneity and is what makes the
result safe. Over both cameras the referent's recolour moves $1.78\times$ as
much image as the placebo's, so the pooled adjusted selectivity is $4.0$. But
the excess is not spread evenly. At the initial state the wrist camera sees the
referent bowl large, its image difference is $2.4\times$ the placebo's, and the
raw $3.0\times$ selectivity falls to $1.3\times$ once that is divided out. At
the later state the two footprints are nearly equal ($1.19\times$) and the
selectivity is $12$ to $16\times$ raw and adjusted alike. So at the state where
footprint could have produced the effect, the adjustment removes most of it;
at the state where it could not, the effect is large and survives untouched.
The adjusted figure exceeds one in every cell, minimum $1.23$, and its cluster
range is $1.27$ to $12.1$. The pixel-count normaliser tells a more conservative
story, with four of eight cells below one: it counts every wrist-camera pixel
of the near bowl as a full unit of perturbation whether it moved by a shade or
a saturation, which is why the norm was pre-registered as primary. Selectivity
holds in all $18$ layers
(Table~\ref{tab:cf-layers}), highest in the first three ($10.8$ to $11.3$) and
lowest in the deep layers that dominate the norm ($5.3$ to $6.2$ at layers $14$
to $17$); the geometric mean over layers is $7.3$, so the conclusion does not
rest on the largest-magnitude layers.

The action follows the features. Recolouring the referent changes the predicted
action chunk by $54\%$ in relative $\ell_2$ pooled, by $12$ to $23\%$ at the
initial state and $89$ to $92\%$ at the later one; recolouring the non-referent
changes it by $2$ to $5\%$ in every cell (Table~\ref{tab:cf-cells}). This is the behavioural form
of H1 and it is expected: the prompt names ``the black bowl'', so a red bowl
is no longer what the noun denotes, and a policy anchored on the referent
must re-plan. The response is sublinear in dose for the target, with elasticity
$0.62$ and near $90\%$ of its full-dose value at half dose, which is consistent
with a categorical rather than graded reading of the colour; the placebo's
response is near linear ($1.21$). Against the positive control, the referent
recolour produces half the latent response of a prompt swap that changes the
action by $73\%$, and the placebo one fourteenth of it.

\subsubsection{Referent versus position}

% Filled from decision_metrics.json -> h2 (per_prompt, grounding, referent_check).
\begin{table}[t]
\caption{Prompt-swap control over identical images and noise. Selectivity is
within-prompt. The anchor is Eq.~\ref{eq:anchor}.}
\label{tab:cf-prompt}
\begin{center}
\begin{tabular}{lcccc}
\hline
Prompt & Target $D$ & Placebo $D$ & Selectivity & Anchor \\
\hline
Task, names the referent      & $10.11$ & $1.42$ & $7.13$ & $17.0$ \\
Sibling, names the placebo    & $10.38$ & $1.75$ & $5.92$ & $9.7$ \\
\hline
\multicolumn{5}{p{5.1in}}{Grounding $g=0.729$ (range $0.680$ to $0.797$ over four (state, draw) pairs); referent moved: no} \\
\hline
\end{tabular}
\end{center}
\end{table}

The swap was processed: it moves the unperturbed action by $73\%$ in relative
$\ell_2$, more than the recolour does. It did not move the referent. Under the
sibling prompt the action remains ten times more sensitive to the original
bowl's colour than to the placebo's (Table~\ref{tab:cf-prompt}), so the policy
is still going for the bowl the task prompt named. The reason is visible in the
scene: the sibling phrase ``next to the ramekin'' is true of a bowl that sits
between the plate and the ramekin. Selectivity weakens from $7.1$ to $5.9$
under the swap, which a rule without the anchor check would have read as
partial evidence that language modulates the response; with the referent
unmoved, that reading is not licensed, and H2 is untestable on this scene with
this prompt. What the swap does show is that a large behavioural effect of
language is compatible with the referent staying put, which bears on how
steering results should be read.

\subsubsection{Circuit corroboration}

% Filled from validation.json (conditions, selectivity, specificity, strata) and nominated_targets.json.
\begin{table}[t]
\caption{Pilot, one traced circuit. Circuit audit for the nominated target $\mathtt{L5}\!:\!0.1\!:\!\mathtt{F735}$
($z^{\mathrm{std}}=+4.86$, exact selectivity $58\times$, active in $39\%$ of
cells). Top: response enrichment of the full circuit, $1681$ parents of which
$99.7\%$ are exercised, over $2000$ matched draws. Bottom: selectivity
enrichment and specificity by influence rank; the head row decides.}
\label{tab:cf-validation}
\begin{center}
\small
\begin{tabular}{lcccc}
\hline
Perturbation & Circuit mean & Random mean & Enrichment & $p$ \\
\hline
Referent recolour     & $0.354$ & $0.275$ & $1.29$ & $0.0005$ \\
Non-referent recolour & $0.035$ & $0.029$ & $1.20$ & $0.0005$ \\
Prompt swap           & $0.685$ & $0.544$ & $1.26$ & $0.0005$ \\
\hline
\end{tabular}

\medskip
\begin{tabular}{lccccc}
\hline
Parents by influence & Target enrich. & $\mathrm{Sel}_{\mathcal{C}}$ & $\mathrm{SelEnrich}$ & $p$ & $\mathrm{Spec}$ \\
\hline
Top 10 (decides) & $1.19$ & $8.8\times$  & $0.92$ & $0.57$  & $1.11$ \\
Top 50           & $1.28$ & $10.5\times$ & $1.11$ & $0.17$  & $0.98$ \\
Top 200          & $1.32$ & $9.8\times$  & $1.04$ & $0.25$  & $0.95$ \\
All ($1676$)     & $1.29$ & $10.0\times$ & $1.07$ & $0.004$ & $1.02$ \\
\hline
\end{tabular}
\end{center}
\end{table}

The nominated target is strongly object-selective on the scene: it responds
$0.344$ to the referent's recolour and $0.006$ to the identical non-referent's,
and fires in $39\%$ of cells. The circuit traced upstream of it has $1681$
parents after pruning, nearly all exercised by the scene. As a set they respond
more than matched random features to every manipulation alike, $1.29\times$
for the referent, $1.20\times$ for the non-referent, $1.26\times$ for the
prompt swap, each at the Monte-Carlo floor (Table~\ref{tab:cf-validation}).
Their target-over-placebo selectivity is $10.0\times$; a random set drawn from
the same layers and flow time gives $9.4\times$, which is the selectivity of
layers $0$ to $4$ themselves (Table~\ref{tab:cf-layers}) weighted by where the
parents sit. The selectivity enrichment is $1.07$, significant over $1676$
nodes and inconsequential. At the head of the influence ranking the top $10$
parents are less selective than random, $0.92$ at $p=0.57$, and the profile is
flat through the top $50$, top $200$ and all: $1.11$, $1.04$, $1.07$. H3 is
falsified for this target. The parents carry no more information about the
perturbed object than any active feature at their layers, and the influence
ranking does not sort by that property at any depth, so a tighter graph would
not change the reading.

Three facts fix what the negative means. It is not about the target, which is
the most object-selective feature the probe found; the failure is in the edges.
It is not the case of parents the scene never exercises, since $99.7\%$ are
active here. And it is invisible to a response-only test: $1.29\times$ at the
floor would read as corroboration under a rule that asks only whether the
parents respond, which is the kind of control the ablations of Section~C are.
What the parents share is responsiveness. Attribution is activation times
gradient, and features that are active and high-variance on the tracing
observations respond more to any change of input. Load-bearing and
content-carrying came apart here, measured.

Two properties of the trace belong on the record. Forty-five of the sixty
expansions were layer-$1$ nodes: with sources drawn from every earlier layer
and a frontier ordered by influence, the layer with the most edges absorbs the
expansion budget, and the result is a broad slice of layers $0$ to $4$ rather
than a layered path. And the target sits at the final flow time, where the
early layers of the action expert process nearly denoised action tokens;
whether the object selectivity those layers show in aggregate is routed to the
target through edges a first-order attribution can see is precisely what this
test says it is not.

\subsubsection{Pooled over tasks}

% Filled from aggregate/aggregate.json -> summary (pooled, bootstrap_ci_95,
% tasks_above_one, null_floor_violations, verdict, h2_verdict) and per_task.
\begin{table}[t]
\caption{Reported run, pooled over the suite's ten tasks at two seeds. Each
pooled figure is a ratio of means over all cells, not a mean of per-cell
ratios, so a task with more usable cells carries more weight. The interval is a
cluster bootstrap that resamples whole tasks, because cells within a task share
a scene and are replicates rather than independent measurements. H2 is gated
per task and never averaged: a task counts as testable only when the anchor of
Eq.~\ref{eq:anchor} exceeds one under the task prompt and falls below one under
the sibling, and a task that fails that gate is untestable, not negative.}
\label{tab:cf-pooled}
\begin{center}
\small
\begin{tabular}{lccc}
\hline
Quantity & Pooled & Task-cluster 95\% CI & Tasks above $1$ \\
\hline
$\mathrm{Sel}$, raw                          & \RESULT{val} & [\RESULT{lo}, \RESULT{hi}] & \RESULT{k}/10 \\
$\mathrm{Sel}_{\mathrm{adj}}$, $\lVert\Delta I\rVert_2$ (decides) & \RESULT{val} & [\RESULT{lo}, \RESULT{hi}] & \RESULT{k}/10 \\
$\mathrm{Sel}_{\mathrm{adj}}$, $|P|$          & \RESULT{val} & [\RESULT{lo}, \RESULT{hi}] & \RESULT{k}/10 \\
Action selectivity $A_t/A_p$                  & \RESULT{val} & [\RESULT{lo}, \RESULT{hi}] & \RESULT{k}/10 \\
\hline
Null floor, worst task                        & \RESULT{val} & --- & --- \\
Cells contributing                            & \RESULT{n}   & --- & --- \\
\hline
\multicolumn{4}{l}{\emph{H2, gated per task}} \\
Referent moved, so H2 is testable  & \RESULT{k}/10 & --- & --- \\
Of those, selectivity inverted      & \RESULT{k}/\RESULT{m} & --- & --- \\
\hline
\end{tabular}
\end{center}
\end{table}

\subsubsection{What the audit can detect}

% Filled from h3/calibration_*/audit_calibration.json -> levels, summary.
\begin{table}[t]
\caption{Audit sensitivity, measured on the pilot's own delta store: sets of
$K=16$ drawn from the $25{,}032$ features exercised below layer $5$ at
$\tau=0.1$, with a fraction $c$ taken from the top of the measured selectivity
ranking (Eq.~\ref{eq:contamination}) and the remainder at random, $R=20$
repeats and $M=2000$ matched draws per repeat. A level counts as detected when
the selectivity enrichment exceeds one at $p<0.05$.}
\label{tab:cf-calibration}
\begin{center}
\small
\begin{tabular}{lccc}
\hline
Content $c$ & Selective members & $\mathrm{SelEnrich}$, median & Detected \\
\hline
$0$            & $0$  & $0.95$   & $5\%$   \\
$1/8$          & $2$  & $1.46$   & $40\%$  \\
$1/4$          & $4$  & $2.05$   & $85\%$  \\
$1/2$          & $8$  & $3.42$   & $100\%$ \\
$3/4$          & $12$ & $7.37$   & $100\%$ \\
$1$            & $16$ & unbounded & $100\%$ \\
\hline
\multicolumn{4}{l}{Detection floor $c^{\star}=1/4$; same floor at $K=10$ and $K=50$} \\
\hline
\end{tabular}
\end{center}
\end{table}

The audit is neither blind nor over-calling. A fully content-carrying set is
always detected, and at $c=1$ the set has no placebo response at all, so its
selectivity is unbounded rather than merely large. A purely random set is
called $5\%$ of the time, which is the nominal level and not above it. Between
those, detection rises monotonically and crosses the $80\%$ bar at $c=1/4$:
four genuinely selective members out of sixteen are enough. The floor is
insensitive to set size, at $80\%$ for $K=10$ and $95\%$ for $K=50$ at the same
$c$, so it is a property of the statistic rather than of the circuit's size.

This is what makes the pilot's negative say something. The traced circuit was
not merely unenriched: it is less than a quarter content-carrying, since a
circuit of which a quarter were selective would have been detected at $85\%$
and this one was read at $0.92$ with $p=0.57$. The same calibration is run for
every target layer the reported run traces, and a pool on which even $c=1$ goes
undetected voids the verdicts drawn from it rather than supporting them.

\subsubsection{H3 as a rate}

% Filled from h3/h3_rate.json -> summary (corroborated, audited, rate,
% clopper_pearson_95, detection_floor, verdict, reading) and audits.
\begin{table}[t]
\caption{Reported run, twenty circuits traced from targets nominated across the
suite's ten tasks with sources restricted to the previous layer. Each row is
one target; the head is the top ten parents by traced influence, where the
Monte-Carlo $p$ is still informative. The rate at the foot is the decision
quantity, with an exact binomial interval.}
\label{tab:cf-rate}
\begin{center}
\small
\begin{tabular}{lccccc}
\hline
Target & Feature $\mathrm{sel}$ & Parents & Head $\mathrm{SelEnrich}$ & $p$ & $\mathrm{Spec}$ \\
\hline
\RESULT{key} & \RESULT{val} & \RESULT{n} & \RESULT{val} & \RESULT{p} & \RESULT{val} \\
\multicolumn{6}{c}{\emph{$\cdots$ eighteen further rows $\cdots$}} \\
\RESULT{key} & \RESULT{val} & \RESULT{n} & \RESULT{val} & \RESULT{p} & \RESULT{val} \\
\hline
\multicolumn{6}{l}{Corroborated: \RESULT{k} of $20$, Clopper-Pearson 95\% CI [\RESULT{lo}, \RESULT{hi}]} \\
\multicolumn{6}{l}{Detection floor over the pools audited: \RESULT{val}} \\
\multicolumn{6}{l}{Verdict: \RESULT{supported / falsified / inconclusive / invalid}} \\
\hline
\end{tabular}
\end{center}
\end{table}

\subsubsection{Verdicts}
\label{sec:verdicts}

The rules were fixed before the run. None of the thresholds is tuned: the null
bound of $5\%$ is the level below which the floor could not change the
direction of a ratio near one; the grounding bound of $0.01$ is one per cent of
the action norm; the consistency bound of $25\%$ is the weakest rule that
excludes a feature seen under one state or one draw; the depth bound leaves
four layers of parents to trace; the exact-selectivity bound of two is the
smallest integer ratio that makes ``responds more to the referent''
unambiguous; the selectivity and specificity bounds of one are where a circuit
is indistinguishable from a matched random set and where an intervention is no
more informative than an unrelated one; the head of ten is the smallest stratum
at which a mean is still a mean and $p$ still informative; $\alpha=0.05$ is
conventional.

\paragraph{H1.} Supported. The rule is: invalid if the null floor exceeds
$5\%$ of the placebo response; supported if the pooled
$\mathrm{Sel}_{\mathrm{adj}}$ with $S=\lVert\Delta I\rVert_2$ and its minimum
over cells both exceed one; weak if only the pooled value does; falsified
otherwise. The floor is zero, the pooled adjusted selectivity is $4.0$, and its
minimum over the eight cells is $1.23$. The relative-$\ell_2$ ratio ($7.1$) and
the per-layer geometric mean ($7.3$) agree in direction with the pooled figure,
so the conclusion is robust across summaries.

\paragraph{H2.} Untestable. The rule is: untestable if there is no sibling
prompt, if $g<0.01$, or if the anchor of Eq.~\ref{eq:anchor} does not exceed
one under the task prompt and fall below one under the sibling prompt;
otherwise supported if the sibling prompt's selectivity falls below one,
falsified if it is unchanged or larger, partial in between. Here $g=0.73$ but
the anchor is $17$ under the task prompt and $9.7$ under the sibling prompt:
the referent did not move.

\paragraph{H3.} The rule has two levels. Per circuit: inconclusive if fewer
than half the parents are exercised; corroborated if the selectivity enrichment
of the top $10$ parents by influence exceeds one at $p<0.05$ and the full
circuit agrees in direction; not corroborated otherwise; qualified as
object-specific if specificity also exceeds one at the head. Over circuits:
invalid if any audited pool failed its calibration, so that a flat reading
carries no information; supported if the Clopper-Pearson interval on $k/N$
excludes $5\%$; falsified if its upper end falls below $50\%$; inconclusive
otherwise.

For the pilot's single circuit the per-circuit reading is \emph{not
corroborated}: $99.7\%$ of parents are exercised, the head is $0.92$ at
$p=0.57$, and the full circuit is $1.07$. Its calibration passes, so the
reading stands, and Table~\ref{tab:cf-calibration} sharpens it: the circuit is
less than a quarter content-carrying. At $N=1$ the rate is untestable, because
Eq.~\ref{eq:cp-zero} gives a bound of $95\%$ and excludes nothing. The reported
run's verdict is \RESULT{supported / falsified / inconclusive / invalid} at
$k=\RESULT{k}$ of $N=20$.

\subsection{Relation to the rest of the paper}
\label{sec:relation}

\begin{table}[t]
\caption{How this section bears on each preceding one.}
\label{tab:cf-relation}
\begin{center}
\small
\begin{tabular}{p{0.9in}p{1.55in}p{1.6in}p{0.7in}}
\hline
Section & It establishes & This section adds & Effect \\
\hline
Architecture and training & the sparse codes reconstruct the frozen MLPs (held-out loss $0.046$) & the codes are organised by a controlled visual property: $7\times$ pooled, $58\times$ for single features, in all $18$ layers & supports \\
A. Feature discovery & nomination by activation statistics; content by inspection of top examples & nomination by controlled selectivity, and an intervention test of a feature's claimed content & complements \\
B. Circuit construction & attribution-traced compact circuit as the main hypothesis & for a $58\times$-selective target the parents carry no object content beyond their layers; the frontier collapses into the highest-degree layer & challenges, conditionally \\
C. Causal validation & circuit ablation moves the action $2\times$ more than random circuits; target inactive on negatives & load-bearing and content-carrying separated by measurement; a non-tautological negative; a positive control for scale & qualifies \\
Steering & a speed feature steers open-loop actions & object-anchored targets with a $54\%$ behavioural effect; a language manipulation can move the action without moving the referent & supports, extends \\
\hline
\end{tabular}
\end{center}
\end{table}

Table~\ref{tab:cf-relation} summarises the bearing of this section on each
preceding one; the paragraphs below give the reasoning.

\paragraph{Architecture and training.}
Reconstruction loss shows that the transcoders reproduce the MLP outputs; it
does not show that their sparse basis is organised along anything a person
would recognise. H1 is the first evidence in the paper that it is: a
controlled change to one object's appearance moves the code seven times more
than the same change to an identical object the task does not name, in every
layer, and single features reach $58\times$. This supports the use of the
transcoders as a lens on the policy, which every later section assumes.

\paragraph{Feature discovery.}
Section~A nominates by how unusually a feature activates and reads its content
from its top-activating examples. Both are observational, and the section
itself notes that a top feature can be task- or episode-specific rather than a
mechanism. The probe nominates by a different route, controlled selectivity
with a known cause, and its filters reject $4310$ candidates that a single
large observation would otherwise have promoted. It also supplies the test that
inspection cannot: the black-bowl content attributed to F9970 by its top
examples is exactly the kind of claim that a bowl recolour decides. Nothing
here contradicts Section~A's rankings; it gives them a second, independent
nomination and the means to check what they nominate.

\paragraph{Circuit construction.}
This is the section our result presses on, and the pressure is conditional.
The F9970 circuit was pruned to sixteen nodes and shown to move the action; the
circuit we traced with the same tracer for a target at least as selective was
pruned to $1681$ nodes and shown to carry nothing about the object beyond its
layers' baseline, with an influence ranking that does not sort by the property
at any depth. Neither result transfers to the other automatically: the targets,
tasks, pruning and flow times differ. What does transfer is the presumption.
Section~B treats a traced graph as ``the main circuit hypothesis'' for the
target's content; our result shows that this tracer can return a graph that is
a slice of active features in the target's layers, and that the graph's
attribution ranking can be uninformative about the very property the target is
selective for. A circuit hypothesis from this tracer therefore needs a content
test before it is read semantically, and the F9970 circuit has not had one. We
also observe a concrete mechanism, the frontier's collapse into layer $1$, that
Section~B's tighter pruning may have avoided or may have hidden.

\paragraph{Open-loop causal validation.}
Section~C establishes that ablating the traced circuit moves the action twice
as much as ablating matched random circuits, and that the target is inactive on
its negatives. Our result qualifies rather than contradicts it. First, its
random-circuit control is a response-type test: it asks whether the traced set
moves the action more than a random set does. We show a traced set that passes
the analogous input-side test at $1.29\times$ while carrying no content about
the property its target encodes, so the $2\times$ of Section~C establishes
load-bearing and not content-carrying, and the two should be named separately.
Second, Section~C's own observation that parent ablation affects the negatives
at about half its effect on the positives, which it attributes to the parents
being ``broader'', is what a set selected for responsiveness would produce.
Third, its negative sets are inactive by construction; the matched placebo is
the non-tautological control, since the feature is active on the scene and the
pixels change by nearly the same amount. Fourth, the positive control gives the
scale Section~C lacks, with one caveat: Section~C reports the $\ell_2$ of the
chunk difference in action units, this section the same quantity relative to
the chunk norm, and the two are compared only after that conversion.

\paragraph{Steering.}
The steering section moves open-loop actions along a speed axis whose positives
come from one fast trajectory. Two things here support and extend it. The
nominated features are anchored on the task's referent with exact
selectivities from $36\times$ to $130\times$, and recolouring that referent
moves the action by $54\%$, so they are candidates for steering along a
referent axis defined by a controlled visual property rather than by episode
identity. And the prompt-swap result is a caution the steering claims should
carry: a language manipulation moved the action by $73\%$ without moving the
referent at all, so a behavioural check of what an intervention changed, not
only how much, belongs next to every steering effect.

\subsection{Threats to validity}
\label{sec:threats}

Within a pair the manipulation is isolated: state, camera and noise are shared
and the pixel difference is confined to the object. Across the pair the target
and placebo differ in position as well as referent status; the footprint
adjustment addresses the first and the prompt swap was to address the second,
and on this scene it could not. The circuit is traced on dataset frames and
audited on live renders of one task; a parent that is upstream on the tracing
distribution may not respond here, and while the exercised fraction rules out
the case of parents the scene never activates, an exercised parent that does
not respond remains ambiguous between not carrying the property and carrying it
elsewhere. Specificity has two readings, generic responsiveness or a referent
circuit that integrates language, which only selectivity enrichment separates;
in the reported run both point the same way. The traced graph's size is a
pruning setting and a large graph sits near the random mean by construction;
the flat influence profile is what shows the negative here is about the
attribution and not the pruning. All claims are about transcoder features, not
the MLP basis. Colour is load-bearing in this suite because the prompt names
it; in a task whose prompt did not, the same recolour would be expected to move
behaviour little.

Three limits of the reported run bound what the rate can claim. Ten tasks is
the suite's ceiling and not a budget, since a matched placebo needs a duplicated
mesh and only \texttt{libero\_spatial} has one throughout; the interval on H1
rests on ten clusters and will be wide. The twenty audited targets are the
probe's own nominees, so they are the features most likely to have a
content-carrying circuit and the rate is an upper bound on the tracer's
performance rather than an average over features. And the nominations sit near
layer five while the paper's case study is at layer eleven, so a comparison
between the two is confounded with depth until depth-matched nominations are
run.

\subsection{Limitations and open directions}
\label{sec:gaps}

\begin{enumerate}
\item \textbf{Audit the paper's own circuit.} The F9970 compact circuit is the
paper's central example and has been shown load-bearing, not content-carrying.
The same audit applies to it unchanged, on a scene containing the bowl and
drawer, with the bowl recolour as the manipulation. This is the single
experiment that most changes how Sections~B and~C should be read.
\item \textbf{The tracer as a factor.} The reported run restricts sources to
the previous layer, because in the pilot forty-five of sixty expansions
collapsed into layer $1$ and the graph reached $1681$ parents. Running the same
twenty targets under both source policies would separate a property of the
circuits from a property of the frontier. The flat influence profile predicts
the verdict will not change; the prediction is worth recording either way.
\item \textbf{Depth-matched nominations.} The probe nominates near layer five
and the case study's circuit sits at layer eleven. Raising the depth bound
until the nominated targets match it would remove the one confound between the
two results that the audit cannot adjust for.
\item \textbf{An unambiguous sibling prompt.} H2 was untestable because the
sibling phrase was true of both bowls. A prompt override that names only the
placebo's region, with the anchor check as the gate, would make H2 decidable.
\item \textbf{Steering on the referent axis.} The nominated features are the
natural targets for a steering experiment whose positives and negatives are
defined by a controlled visual property, and whose success is read on the
behavioural anchor rather than on a speed score.
\item \textbf{Tracer variants.} The audit distinguishes tracers. Attribution at
the position where the target responds on the audited scene, rather than at
the dataset position; second-order or integrated-gradient edges; and pruning by
content rather than by cumulative influence are all testable against the same
known cause.
\item \textbf{Breadth beyond one suite.} Ten tasks is
\texttt{libero\_spatial}'s ceiling, because a matched placebo needs a duplicated
mesh. \texttt{libero\_10} adds two more through its \texttt{moka\_pot} pairs,
which is the only way to raise the cluster count without weakening the match.
Beyond that the placebo has to be a different object, as it would be on the
drawer task of the case study, and the footprint adjustment carries more of the
weight. A second object type and a non-colour perturbation such as texture
would each widen what the design can claim.
\item \textbf{Units.} Section~C and this section measure the action in
different normalisations of the same quantity; one column in the same unit
would let the positive control calibrate Section~C's effect sizes directly.
\item \textbf{Closed loop.} All measurements are open-loop at a frozen state.
Whether a $54\%$ change in the predicted chunk changes what the arm does
requires rollouts, where task success and the object reached can be measured.
\end{enumerate}

\subsection{Reproducibility}

Every stage is a script with recorded configuration and a deterministic seed;
noise draws are seeded per (state, draw) so any cell can be regenerated alone.
Each run writes its full per-feature deltas and baselines for every condition,
every measurement including the controls, the decision quantities and verdicts
with the rules they were decided under, the tensor shapes handed to the policy,
and a provenance record with the git commit, library versions, device and a
content hash of the transcoder checkpoint. Runs execute on an ephemeral notebook machine, so each probe run is copied
to shared storage the moment it finishes rather than at the end of the loop,
and the copy is verified by byte count because a full quota returns success on
write. Each batch also carries a manifest listing every run, its verdicts, its
commit and where each quantity is read from, so the analysis can be repeated
from the stored directory alone without the session that produced it. The reported run was produced at commit \texttt{758b7c9} of the
repository with a clean tree, on an NVIDIA L4 with Python 3.12, LeRobot 0.6.1,
PyTorch 2.11, MuJoCo 3.8.1 and robosuite 1.4.0, against a transcoder
checkpoint of SHA-256 prefix \texttt{b75ee195bbe11ede}. The pipeline is probe, report, discovery, trace, validate; the
notebook orchestrates these and displays their artefacts without computing any
decision itself. A guard aborts the run rather than continuing on a stale render, an
unrestored baseline, drifting tensor shapes, a policy whose weights did not
fully load, a nomination with no eligible feature, or a prompt swap that did
not move the behavioural referent. A perturbation that changed no pixels skips
that state and is counted, since over ten tasks one occluded object should not
discard nine good ones. Every decision quantity in this section is computed by
a script that carries its own tests and is read back from the artefact it
writes; the notebook displays verdicts and computes none.
